% META: {"id":"TEXP-CTP-ME-006","chapitre":"TEXP-COMPLEXES-TRIGO-POLYNOMES","type_objet":"methode","capacites_codes":["C6"],"methodes":["M6"],"mode_creation":"ex_nihilo","sources_inspiration":[],"status":"approved","origin":"needs_review","status_history":["needs_review","audited","accepted","approved"],"release_acceptance":"RELEASE_OWNER_BATCH_ACCEPTANCE","acceptance_closure_digest":"sha256:9b3ccf9a81c5520fbb7e03b7b2d3e7bf2057a02834b6d908bf3be5f49f3b553f","acceptance_receipt_digest":"sha256:4fad86f0282e0fa4e0419cca1ad6379ae7af5a280c04a4a90880f90a1c044242"}
% BEGIN-VERIFY
% from sympy import I, simplify, im, re, Abs
% A = 1 + 2*I; B = 3 + 3*I; C = 7 + 5*I
% r = simplify((C - A)/(B - A))
% assert im(r) == 0 and r == 3
% A2 = 0; B2 = 2 + 2*I; C2 = -1 + I
% r2 = simplify((C2 - A2)/(B2 - A2))
% assert re(r2) == 0 and r2 == I/2
% assert Abs(B - A) == Abs(3 + 3*I - 1 - 2*I)
% END-VERIFY
\begin{fichemethode}{M6}{Étudier des configurations du plan avec les nombres complexes}

\quandUtiliser{%
  Utiliser cette méthode lorsque l'énoncé demande d'étudier alignement,
  orthogonalite, longueurs, angles ou un ensemble de points, a l'aide des
  affixes.
}

\pasApas{%
  \begin{enumerate}
    \item Traduire les objets : longueur $AB = |z_B - z_A|$ ; angle
      $(\vec{AB}\,;\,\vec{AC}) = \arg\dfrac{z_C - z_A}{z_B - z_A}\ [2\pi]$.
    \item Former le quotient cle $Z = \dfrac{z_C - z_A}{z_B - z_A}$ et le
      mettre sous forme algébrique.
    \item Interpréter : $Z$ RÉEL $\iff$ $A$, $B$, $C$ alignes ;
      $Z$ IMAGINAIRE PUR $\iff$ $(AB) \perp (AC)$ ; $|Z| = \frac{AC}{AB}$
      compare les longueurs.
    \item Ensembles de points : traduire la condition en égalité de modules
      ($|z - z_A| = |z - z_B|$ : mediatrice ; $|z - z_A| = r$ : cercle) ou
      d'arguments (demi-droites), puis identifier l'ensemble.
    \item Conclure GEOMETRIQUEMENT en citant la caractérisation utilisee.
  \end{enumerate}
}

\exempleRedige{%
  1) $A$, $B$, $C$ d'affixes $1 + 2i$, $3 + 3i$, $7 + 5i$ sont-ils alignes ?
  2) $A'$, $B'$, $C'$ d'affixes $0$, $2 + 2i$, $-1 + i$ : montrer
  $(A'B') \perp (A'C')$.

  \medskip
  \commentaireMarge{Quotient réel $=$ alignement.}%
  1) $\dfrac{z_C - z_A}{z_B - z_A} = \dfrac{6 + 3i}{2 + i}
  = \dfrac{3(2 + i)}{2 + i} = 3 \in \mathbb{R}$ : les points sont alignes
  (et $\vec{AC} = 3\,\vec{AB}$).

  \medskip
  \commentaireMarge{Quotient imaginaire pur $=$ angle droit.}%
  2) $\dfrac{z_{C'} - z_{A'}}{z_{B'} - z_{A'}} = \dfrac{-1 + i}{2 + 2i}
  = \dfrac{(-1 + i)(2 - 2i)}{8} = \dfrac{4i}{8} = \dfrac{i}{2}$,
  imaginaire pur non nul : $(A'B') \perp (A'C')$.
}

\pieges{%
  \begin{itemize}
    \item \textbf{Piège 1.} Confondre les criteres : réel $=$ alignement,
      imaginaire pur $=$ orthogonalite — ne pas les intervertir.
    \item \textbf{Piège 2.} Former le quotient avec des differences
      incoherentes : les trois affixes doivent partir du MÊME point $A$.
    \item \textbf{Piège 3.} Oublier le cas $Z = 0$ ou un dénominateur nul
      (points confondus) avant de diviser.
  \end{itemize}
}

\verifier{%
  Contrôler le résultat sur une figure (placer les points), et recouper par une
  autre voie : vecteurs colineaires pour l'alignement, produit scalaire nul
  pour l'orthogonalite.
}

\sEntrainer{\refExos{M6}}

\end{fichemethode}
